% !TEX root = ../algorithms.tex

\section{Задача нахождения минимального остовного дерева. Алгоритмы Ярника--Прима--Дейкстры и Краскала. Системы непересекающихся множеств (Union-Find)}

\begin{Definition}{}{}
	Пусть $G=(V,E,c)$ --- связный неориентированный граф с весовой функцией $c\colon E \to \mathbb{R}$. \emph{Остовом} графа $G$ называется подграф $T \subseteq G$, являющийся деревом (связным подграфом без циклов) и содержащий все вершины множества $V$. Остов $T$ называется \emph{минимальным} (МОД), если его суммарный вес $c(T) = \sum_{e \in E(T)} c(e)$ минимален среди всех остовов графа $G$.
\end{Definition}

\begin{Remark}{}{}
	Размер входа равен $O(|V|+|E|)$. Граф удобно хранить списками смежности: для каждой вершины $v \in V$ хранится список пар $(u, c(v,u))$ для всех рёбер $vu \in E$ (структура вида <<вектор векторов пар>>). Это представление особенно удобно для алгоритмов прим-дейкстровского типа, где на каждом шаге требуется перебрать все рёбра, инцидентные текущей вершине.
\end{Remark}

\begin{figure}[H]
	\centering
	\begin{tikzpicture}[scale=0.85, >=Stealth, font=\small]
		\draw (0,0) rectangle (0.7,2.8);
		\draw (0,0.7) -- (0.7,0.7);
		\draw (0,1.4) -- (0.7,1.4);
		\draw (0,2.1) -- (0.7,2.1);
		\node at (0.35,2.45) {$g[0]$};
		\node at (0.35,1.75) {$g[1]$};
		\node at (0.35,1.05) {$g[2]$};
		\node at (0.35,0.35) {$g[3]$};
		\draw (0.7,2.45) -- (3,2.45);
		\node[right] at (3,2.45) {$(u_1,c)\;(u_2,c)\;\ldots$};
		\draw (0.7,1.75) -- (2.5,1.75);
		\node[right] at (2.5,1.75) {$(u_1,c)$};
		\draw (0.7,1.05) -- (3.2,1.05);
		\node[right] at (3.2,1.05) {$(u_1,c)\;(u_2,c)\;(u_3,c)$};
		\draw (0.7,0.35) -- (2.2,0.35);
		\node[right] at (2.2,0.35) {$(u_1,c)$};
	\end{tikzpicture}
	\caption{Хранение графа списками смежности: для каждой вершины $v$ хранится список пар $(u, c(v,u))$ по всем инцидентным рёбрам $vu$. Доступ ко всем рёбрам, инцидентным вершине, занимает время, пропорциональное её степени.}
\end{figure}

\subsection{Лемма о разрезе}

\begin{Definition}{}{}
	\emph{Разрезом} графа $G=(V,E)$ называется пара непустых множеств $(W, V \setminus W)$, $W \neq \emptyset$, $V \setminus W \neq \emptyset$. Ребро $uv \in E$ называется \emph{ребром разреза}, если один из его концов лежит в $W$, а другой --- в $V \setminus W$.
\end{Definition}

\begin{Lemma}{о минимальном ребре разреза}{}
	Пусть $(W, V \setminus W)$ --- разрез графа $G$ и
	\[
	vw = \operatorname*{arg\,min}_{\substack{v \in V\setminus W,\ w \in W\\ vw \in E}} c(v,w)
	\]
	--- ребро разреза минимального веса. Тогда существует минимальный остов $T$ графа $G$, для которого $vw \in E(T)$.
\end{Lemma}

\begin{figure}[H]
	\centering
	\begin{tikzpicture}[>=Stealth, every node/.style={font=\small}]
		\draw[rounded corners=15pt, thick] (-2.2,-1.3) rectangle (-0.2,1.3);
		\draw[rounded corners=15pt, thick] (0.8,-1.3) rectangle (2.8,1.3);
		\node at (-1.2,1.6) {$W$};
		\node at (1.8,1.6) {$V\setminus W$};
		\node[circle,fill,inner sep=1.3pt] (w) at (-0.6,0.4) {};
		\node[circle,fill,inner sep=1.3pt] (v) at (1.2,0.4) {};
		\node[circle,fill,inner sep=1.3pt] (wp) at (-0.6,-0.6) {};
		\node[circle,fill,inner sep=1.3pt] (vp) at (1.2,-0.6) {};
		\node[left] at (w) {$w$};
		\node[right] at (v) {$v$};
		\node[left] at (wp) {$w'$};
		\node[right] at (vp) {$v'$};
		\draw[thick] (w) -- (v);
		\node[above] at (0.3,0.6) {$vw$};
		\draw[dashed] (wp) -- (vp);
		\node[below] at (0.3,-0.8) {$v'w'$};
	\end{tikzpicture}
	\caption{Лемма о разрезе. $vw$ --- ребро разреза минимального веса (сплошное). При добавлении $vw$ к остову $T'$ возникает цикл, в котором найдётся ещё одно ребро разреза $v'w'$ (пунктирное). Замена $v'w'$ на $vw$ не увеличивает вес остова и даёт МОД, содержащий $vw$.}
\end{figure}

\begin{proof}
	\begin{enumerate}
		\item Пусть $T'$ --- произвольный минимальный остов графа $G$.
		\item Если $vw \in E(T')$, то утверждение верно при $T = T'$.
		\item Пусть $vw \notin E(T')$. Так как $T'$ --- дерево, граф $T' + vw$ содержит ровно один цикл, проходящий через ребро $vw$.
		\item Этот цикл переходит из $V \setminus W$ в $W$, поэтому помимо $vw$ в нём найдётся ещё одно ребро разреза $v'w'$ с $v' \in V \setminus W$, $w' \in W$.
		\item Положим $T = T' + vw - v'w'$. Любой путь в $T'$, проходящий через ребро $v'w'$, можно заменить участком найденного цикла, содержащим $vw$, поэтому $T$ остаётся связным. Кроме того, $|E(T)| = |E(T')| = |V|-1$, значит $T$ --- дерево, то есть остов.
		\item По выбору $vw$ имеем $c(v,w) \leq c(v',w')$, следовательно $c(T) = c(T') + c(v,w) - c(v',w') \leq c(T')$.
		\item Поскольку $T'$ минимален, $c(T) \geq c(T')$; вместе с предыдущим неравенством $c(T) = c(T')$. Значит $T$ --- также минимальный остов, причём $vw \in E(T)$.
	\end{enumerate}
\end{proof}

\subsection{Алгоритм Прима (Ярника--Прима--Дейкстры)}

Алгоритм Прима строит МОД, постепенно расширяя поддерево $T$: на каждом шаге к $T$ присоединяется ребро минимального веса, соединяющее уже построенное дерево с оставшимися вершинами.

\begin{algorithmic}
	\For{$v \in V$}
	\State $d_v \gets \infty$, $\quad \pi_v \gets \mathrm{null}$
	\EndFor
	\State $d_{v_0} \gets 0$ \Comment{$v_0$ --- произвольная стартовая вершина}
	\For{$k = 1, \ldots, |V|-1$}
	\State $v \gets \displaystyle\operatorname*{arg\,min}_{u \in V\setminus V(T)} d_u$
	\State $T \gets T \cup \{(v,\pi_v)\}$
	\For{$vw \in E$}
	\If{$d_w > c(v,w)$}
	\State $d_w \gets c(v,w)$
	\State $\pi_w \gets v$
	\EndIf
	\EndFor
	\EndFor
\end{algorithmic}

\begin{Theorem}{Корректность алгоритма Прима}{}
	На каждом шаге алгоритма ребро $(v, \pi_v)$, добавляемое в $T$, принадлежит некоторому минимальному остову графа $G$, содержащему все рёбра, уже добавленные в $T$ ранее.
\end{Theorem}

\begin{proof}
	Индукция по числу добавленных рёбер.
	
	\emph{База.} В начале $T = \emptyset$ --- подмножество любого минимального остова.
	
	\emph{Шаг.} Пусть $T'$ --- минимальный остов, содержащий все рёбра, добавленные в $T$ к текущему моменту, и пусть $W = V(T)$ --- множество вершин построенного поддерева. Так как $d_u$ для $u \in V\setminus V(T)$ равно минимальному весу ребра, соединяющего $u$ с $T$, выбор $v = \operatorname*{arg\,min}_{u \in V\setminus V(T)} d_u$ означает, что $(v, \pi_v)$ --- ребро минимального веса среди всех рёбер разреза $(W, V \setminus W)$.
	\begin{enumerate}
		\item Если $(v, \pi_v) \in E(T')$, утверждение для шага доказано.
		\item Если $(v, \pi_v) \notin E(T')$, то, поскольку $T'$ связен, в $T'$ найдётся некоторое ребро $(x,x') \in E(T')$ разреза $(W, V \setminus W)$. По лемме о разрезе дерево $T'' = T' - (x,x') + (v,\pi_v)$ также является минимальным остовом и содержит все рёбра $T$.
	\end{enumerate}
	По индукции после $|V|-1$ шагов получаем минимальный остов.
\end{proof}

\begin{Remark}{}{}
	Если $\operatorname*{arg\,min}$ вычисляется наивно (перебором), время работы составляет $O(|V|^2)$ --- это оправдано для плотных (в частности, полных) графов. С очередью с приоритетами время работы составляет $O(|E| \log |V|)$.
\end{Remark}

\subsection{Алгоритм Краскала}

Алгоритм Краскала строит МОД, рассматривая рёбра в порядке возрастания веса и поддерживая лес деревьев: ребро добавляется, если оно соединяет вершины из разных деревьев.

\begin{algorithmic}
	\State Отсортировать рёбра так, что $c(e_1) \leq c(e_2) \leq \ldots \leq c(e_{|E|})$
	\State Инициализировать лес: каждая вершина --- отдельное дерево
	\For{$k = 1,\ldots,|E|$}
	\If{ребро $e_k$ соединяет вершины из разных деревьев леса $T_1$ и $T_2$}
	\State добавить $e_k$ в остов
	\State заменить $T_1, T_2$ на новое дерево $T_1 \cup T_2 \cup \{e_k\}$
	\EndIf
	\EndFor
\end{algorithmic}

\begin{Remark}{Корректность}{}
	Корректность доказывается аналогично корректности алгоритма Прима с использованием леммы о разрезе: если $e_k = (v,w)$ соединяет вершины $v$ и $w$ из разных деревьев $T_1$ и $T_2$ текущего леса, то $e_k$ --- ребро минимального веса среди рёбер разреза $(W, V \setminus W)$, где $W$ --- множество вершин дерева $T_1$ (все рёбра меньшего веса либо уже добавлены, либо соединяют вершины внутри одного дерева и потому не пересекают этот разрез).
\end{Remark}

\subsection{Система непересекающихся множеств (Union-Find)}

Для эффективной проверки <<принадлежат ли $v$ и $w$ разным деревьям леса>> в алгоритме Краскала используется структура данных СНМ (Union-Find).

\begin{Definition}{}{}
	\emph{Система непересекающихся множеств} (СНМ) --- структура данных, поддерживающая разбиение множества элементов $\{x_1,\ldots,x_n\}$ на непересекающиеся подмножества с двумя операциями:
	\begin{enumerate}
		\item $\mathrm{union}(x, x')$ --- объединяет множества, содержащие $x$ и $x'$, в одно новое множество, удаляя старые;
		\item $\mathrm{find}(x)$ --- находит представителя множества, содержащего $x$.
	\end{enumerate}
\end{Definition}

\begin{Statement}{}{}
	Элементы $x$ и $y$ лежат в одном множестве $\iff \mathrm{find}(x) = \mathrm{find}(y)$.
\end{Statement}

\begin{Remark}{}{}
	Каждый элемент $x$ хранит указатель $x.par$ на родителя. Если $x.par = x$, то $x$ --- представитель своего множества. Таким образом множества представляются деревьями (лесом), а $\mathrm{find}(x)$ --- это переход по указателям $par$ до корня.
\end{Remark}

\begin{figure}[H]
	\centering
	\begin{tikzpicture}[>=Stealth, font=\small]
		\node (r) at (0,0) {$r$};
		\node (a) at (-1.2,1.2) {$x_1$};
		\node (b) at (1.2,1.2) {$x_2$};
		\node (c) at (-1.8,2.4) {$x_3$};
		\node (d) at (-0.6,2.4) {$x_4$};
		\node[right] at (0.2,-0.1) {$r.par=r$};
		\draw[->] (a) -- (r);
		\draw[->] (b) -- (r);
		\draw[->] (c) -- (a);
		\draw[->] (d) -- (a);
	\end{tikzpicture}
	\caption{Множество как дерево: указатели $x.par$ ведут к корню $r$ --- представителю множества ($r.par = r$). Операция $\mathrm{find}(x)$ поднимается по указателям от $x$ до корня.}
\end{figure}

\subsubsection*{Эвристика 1 (объединение по высоте)}

Каждый элемент $x$ хранит поле $x.h$ --- высоту дерева с корнем $x$.

\begin{algorithmic}
	\Function{union}{$x, y$}
	\State $x \gets \mathrm{find}(x)$
	\State $y \gets \mathrm{find}(y)$
	\If{$x = y$} \Return \EndIf
	\If{$x.h < y.h$}
	\State $x.par \gets y$
	\Else
	\State $y.par \gets x$
	\If{$x.h = y.h$}
	\State $x.h \gets x.h + 1$
	\EndIf
	\EndIf
	\EndFunction
\end{algorithmic}

\subsubsection*{Эвристика 2 (сжатие путей)}

\begin{algorithmic}
	\Function{find}{$x$}
	\If{$x.par \neq x$}
	\State $x.par \gets \mathrm{find}(x.par)$
	\EndIf
	\State \Return $x.par$
	\EndFunction
\end{algorithmic}

\begin{Lemma}{}{}
	При использовании эвристики 1 для любого элемента $x$ поддерево с корнем $x$ содержит не менее $2^{x.h}$ элементов.
\end{Lemma}

\begin{proof}
	Индукция по $h = x.h$.
	
	\emph{База} ($h = 0$): поддерево с корнем $x$ состоит из одного элемента, $2^0 = 1$ --- верно.
	
	\emph{Переход}: значение $x.h$ увеличивается с $h-1$ до $h$ только в операции $\mathrm{union}$, когда дерево с корнем $x$ высоты $h-1$ объединяется с деревом высоты $h-1$, корень которого становится сыном $x$. По предположению индукции каждое из этих двух поддеревьев содержит не менее $2^{h-1}$ элементов, поэтому поддерево с корнем $x$ содержит не менее $2^{h-1} + 2^{h-1} = 2^h$ элементов.
\end{proof}

\begin{figure}[H]
	\centering
	\begin{tikzpicture}[>=Stealth, font=\small]
		\coordinate (root) at (0,2);
		\coordinate (l) at (-1,0.6);
		\coordinate (rr) at (1,0.6);
		\draw (root) -- (l);
		\draw (root) -- (rr);
		\draw (l) +(-0.7,-0.7) -- (l) -- +(0.7,-0.7) -- cycle;
		\draw (rr) +(-0.7,-0.7) -- (rr) -- +(0.7,-0.7) -- cycle;
		\node[above] at (root) {высота $h$};
		\node at (l) {$\ge 2^{h-1}$};
		\node at (rr) {$\ge 2^{h-1}$};
	\end{tikzpicture}
	\caption{Дерево высоты $h$ образуется объединением двух деревьев высоты $h-1$; по индукции каждое содержит $\ge 2^{h-1}$ элементов, а вместе --- $\ge 2^h$. Отсюда высота ограничена логарифмом числа элементов.}
\end{figure}

\begin{Theorem}{}{}
	При использовании эвристики 1 высота любого дерева в СНМ не превышает $\log_2 n$, где $n$ --- общее число элементов. Следовательно, операции $\mathrm{find}$ и $\mathrm{union}$ выполняются за время $O(\log n)$.
\end{Theorem}

\begin{proof}
	Высота дерева с корнем $x$ не превышает $x.h$. По предыдущей лемме поддерево с корнем $x$ содержит не менее $2^{x.h}$ элементов, а так как всего элементов не более $n$, имеем $2^{x.h} \leq n$, откуда $x.h \leq \log_2 n$. Значит, высота любого дерева не превышает $\log_2 n$, и обход от произвольного элемента до корня в операциях $\mathrm{find}$ и $\mathrm{union}$ занимает $O(\log n)$ шагов.
\end{proof}

\subsection{Оценка сложности алгоритма Краскала}

\begin{Remark}{}{}
	Сортировка рёбер занимает время $O(|E| \log |E|)$, каждая операция с СНМ --- $O(\log |V|)$, итого $O(|E| \log |E|)$. Так как граф простой, $|E| \leq |V|^2$, а значит $\log |E| = O(\log |V|)$. Поэтому итоговое время работы алгоритма Краскала
	\[
	O(|E| \log |E|) = O(|E| \log |V|).
	\]
\end{Remark}

\begin{Remark}{}{}
	С обеими эвристиками (объединение по высоте и сжатие путей) последовательность операций СНМ выполняется за $O(|E| \cdot \alpha(|V|))$, где $\alpha$ --- обратная функция Аккермана (растёт крайне медленно). Однако сортировка рёбер всё равно остаётся узким местом, поэтому общая оценка алгоритма Краскала не улучшается.
\end{Remark}